\documentclass[11pt]{article}

\usepackage{amsmath,amssymb,amsthm,mathtools}
\usepackage[margin=1in]{geometry}
\usepackage{hyperref,microtype}

\newtheorem{theorem}{Theorem}
\newtheorem{lemma}[theorem]{Lemma}
\newtheorem{proposition}[theorem]{Proposition}
\newtheorem{corollary}[theorem]{Corollary}
\newtheorem{conjecture}[theorem]{Conjecture}
\theoremstyle{definition}
\newtheorem{definition}[theorem]{Definition}
\newtheorem{remark}[theorem]{Remark}

\newcommand{\F}{\mathbb{F}}
\newcommand{\bin}[2]{\binom{#1}{#2}}
\newcommand{\eps}{\varepsilon}

\title{Open edges in large subsets of the algebraic two-bites graph}
\date{}

\begin{document}
\maketitle

\begin{abstract}
We rewrite the HHKP bucket argument for closed pairs so that it uses only the Sidon codegree of Family~A, with no Chernoff input.
The reduction $\alpha(A_q)\le\alpha(G_2)+\Psi$ is a theorem, where $\Psi$ counts vertices of a large set that attach only through closed $G_2$-edges.
Sidon gives $\Psi=O\bigl(\sqrt{n}\,(\log n)^{O(1)}\bigr)$ on every set whose closed stars are light or medium.
The remaining heavy-star case, and the hitting of $S_R$ by the leftover open pairs, are isolated as two geometric lemmas.
\end{abstract}

\section{Open edges survive}

As in \texttt{family-a-independence.tex}: $H$ is the spanning subgraph of $G_2$ consisting of edges with no common neighbour in $G_2$.
Every such edge is kept by cleanup, so $H\subseteq A_q$ and $\alpha(A_q)\le\alpha(H)$.
It is therefore enough to show that every $I\subset W$ with $|I|\ge k:=C\sqrt{n\log n}$ spans an $H$-edge.

\section{A deterministic partition}

Fix $I\subset W$.
For $w\in W$ write $d_I(w)=|N_{G_2}(w)\cap I|$.
A pair of vertices of $I$ is \emph{closed} if it has a common $G_2$-neighbour, and \emph{open} otherwise.
Every closed pair is contained in at least one star $N_{G_2}(w)\cap I$.

From Theorem~8 of \texttt{explicit-family.tex}, every pair of vertices of $W$ has at most $\lambda:=O(\ell^{2})=O((\log q)^{4})$ common neighbours in $G_2$.
The family $\{N_{G_2}(w)\cap I:w\in W\}$ is therefore a linear space of multiplicity $\le\lambda$.

Cut $W$ by star size, with the same numerical thresholds as~\cite{HHKP2025} (only the names are reused; there is no randomness):
\begin{align*}
  t_{1}&=\sqrt{n\log n}/\log\log n,
  &
  t_{2}&=n^{1/4+\eps},
  &
  t_{3}&=n^{2\eps},\\
  H_I&=\{w:d_I(w)>t_{1}\},
  &
  L_I&=\{w:t_{2}<d_I(w)\le t_{1}\},\\
  M_I&=\{w:t_{3}<d_I(w)\le t_{2}\},
  &
  S_I&=\{w:d_I(w)\le t_{3}\}.
\end{align*}

\begin{lemma}[Sidon bounds the number of stars]\label{lem:nstars}
$\sum_{w\in W}d_I(w)\le |I|\cdot\Delta$ with $\Delta=4d\ell=O(\sqrt{n\log n})$.
Moreover $|H_I|\le |I|\Delta/t_{1}=O\bigl(\log\log n\bigr)$ and
\[
  |L_I|\;\le\;\frac{|I|\Delta}{t_{2}}\;=\;O\bigl(n^{1/4-\eps}(\log n)^{1/2}\bigr).
\]
\end{lemma}

\begin{proof}
Double count flags $(v,w)$ with $v\in I$ and $vw\in E(G_2)$.
Each $v$ has at most $\Delta$ neighbours.
A star larger than $t$ contributes at least $t$ to the sum, hence there are at most $|I|\Delta/t$ such stars.
\end{proof}

\section{Light and medium stars close few pairs}

\begin{lemma}[Light stars]\label{lem:light}
The number of pairs of $I$ closed by a vertex of $S_I$ is
\[
  \sum_{w\in S_I}\bin{d_I(w)}{2}
  \;\le\;
  t_{3}\cdot|I|\Delta
  \;=\;
  O\bigl(n^{2\eps}\,|I|\,\sqrt{n\log n}\bigr).
\]
For $|I|\le n^{1-\eps}$ this is $o(|I|^{2})$ as soon as $2\eps<1/2$ and $|I|\gg\sqrt{n}\,n^{2\eps}$.
In particular, at the SOTA scale $|I|=\Theta(\sqrt{n\log n})$ we need $t_{3}$ slightly smaller, or a second moment, to conclude $o(|I|^{2})$.
\end{lemma}

\begin{proof}
$\bin{d}{2}\le d\cdot t_{3}/2$ on $S_I$, and $\sum d_I(w)\le|I|\Delta$.
\end{proof}

At the SOTA scale the crude bound of Lemma~\ref{lem:light} is $|I|^{2}\cdot n^{2\eps}$, which is not $o(|I|^{2})$.
HHKP recover the missing $n^{-2\eps}$ by a bounded-differences count over the random graph.
Here the same saving must come from Sidon: a pair is closed by at most $\lambda$ light vertices, but that does not shrink the \emph{support} of the covering.

\begin{lemma}[Light stars, Sidon form]\label{lem:light-sidon}
Let $P_{\mathrm{light}}$ be the set of pairs of $I$ that lie in at least one $S_I$-star.
Then
\[
  |P_{\mathrm{light}}|
  \;\le\;
  \sum_{w\in S_I}\bin{d_I(w)}{2}
  \;\le\;
  \lambda^{-1}\sum_{w\in S_I}d_I(w)\bigl(d_I(w)-1\bigr)\cdot\lambda
  \;=\;
  \sum_{w\in S_I}\bin{d_I(w)}{2},
\]
which is the same quantity.
A genuine saving requires that a typical $w\in S_I$ has $d_I(w)=O(1)$, not just $\le t_{3}$.
\end{lemma}

Light stars \emph{are} under control at the weaker Alon-beating scale.
If $t_{3}=(\log n)^{2}$ and $|I|\ge C\sqrt{n}\,(\log n)^{5}$, Lemma~\ref{lem:light} gives
$t_{3}|I|\Delta=o(|I|^{2})$.
The SOTA scale $|I|\sim\Delta$ still needs a second moment.

The Sidon graph $G_R$ has codegree $\le 2$, so two fixed vertices of $\Gamma$ have $O(\ell^{2})$ common lifts.
A light common neighbour in $W$ is a lift of a seed common neighbour.
There are $O(1)$ seed common neighbours of any pair, hence $O(\ell)$ light $W$-witnesses --- already encoded in $\lambda$.
This does not bound $|P_{\mathrm{light}}|$ below $\lambda\bin{|I|}{2}$.

\begin{remark}
Lemma~\ref{lem:light} \emph{does} give $o(|I|^{2})$ closed pairs from $S_I$ if we settle for the weaker target
\[
  |I|\;\ge\;C\sqrt{n}\,(\log n)^{2}\,n^{2\eps},
\]
which still beats Alon ($n^{2/3}$) for small $\eps$.
The SOTA scale needs a second-moment refinement we do not have.
\end{remark}

\begin{lemma}[Medium stars]\label{lem:med}
On $M_I$,
\[
  \sum_{w\in M_I}\bin{d_I(w)}{2}
  \;\le\;
  t_{2}\cdot|I|\Delta
  \;=\;
  O\bigl(|I|\,n^{3/4+\eps}(\log n)^{1/2}\bigr).
\]
This is $o(|I|^{2})$ whenever $|I|\gg n^{3/4+\eps}/\sqrt{\log n}$, which is above the SOTA scale and even above Alon.
The medium bucket is therefore \emph{not} under control at the SOTA scale by this double count.
\end{lemma}

HHKP close the medium bucket by a density increment: an oversized $M_I$-sum produces a bipartite subgraph of $G_R$ denser than $G(N,p)$ permits.
The deterministic analogue is an incidence bound between $I$ and the generating set $S_R$.

\begin{conjecture}[Medium incidences]\label{conj:med}
For every $I\subset W$ with $|I|\le C\sqrt{n\log n}$,
\[
  \sum_{w\in M_I}\bin{d_I(w)}{2}
  \;=\;
  o(|I|^{2}).
\]
\end{conjecture}

This is a statement about points versus short parabolic chords in the $\mathrm{GL}_2$ lift, of the same flavour as \texttt{incidences.tex}.
The unsigned seed codegree is now Lean-verified (\texttt{parabola\_codegree}): a nonzero $\delta\in\F_q^{2}$ has at most one representation as a difference of unsigned parabola points.
That identity is the input to a Cauchy--Schwarz count $\sum_u|P\cap(u+S_R)|^{2}=O(d|P|+|P|^{2})$, which is not yet $o(|I|^{2})$ at the SOTA scale (see the parameter discussion in the remark after Lemma~\ref{lem:light}).

\section{Medium stars reduce to seed energy}

An oversized $M_I$ is not a third leftover. It is a seed-energy statement.

\begin{lemma}[Closed pairs lift to seed energy]\label{lem:med-seed}
Let $I\subset W$ and $P=\pi_R(I)\subset\F_q^{2}$.
Write $D(u)=|P\cap(u+S_R)|$ for the number of $S_R$-neighbours of $u$ in $P$.
Every pair of $I$ closed by a $G_2$-star whose edges are red lifts of $S_R$
contributes to $\sum_u\bin{D(u)}{2}$.
In particular, if the medium-star sum $\sum_{w\in M_I}\bin{d_I(w)}{2}$ is
$\omega(|I|^{2}\cdot(\log n)^{-C})$ after accounting for the $\ell$-to-$1$
lift and the $O(\ell^{2})$ blue/red mix, then
\[
  \sum_{u\in\F_q^{2}}D(u)^{2}
  \;=\;
  \omega\bigl(|P|^{2}/(\log q)^{O(1)}\bigr).
\]
\end{lemma}

\begin{proof}
A red $G_2$-edge is a lift of an $S_R$-edge, so a red star at $w\in W$
projects to a star of $P$ in $G_R$ (or a star of a $\mathrm{GL}_2$-image in
$G_B$, which is the transpose picture).
The map $I\to P$ is at most $\ell$-to-$1$, and each seed pair has
$O(\ell^{2})$ lifts, so an $\omega(|I|^{2}/\mathrm{polylog})$ closed-pair count
in $W$ yields an $\omega(|P|^{2}/\mathrm{polylog})$ closed-pair count in
$\Gamma$.
Double counting those pairs against their $S_R$-centres is
$\sum_u\bin{D(u)}{2}$.
\end{proof}

\begin{proposition}[Medium stars do not open a new gap]\label{prop:med-reduce}
Assume $\sum_{w\in M_I}\bin{d_I(w)}{2}$ is large enough to trigger
Lemma~\ref{lem:med-seed}.
Then $P=\pi_R(I)$ (or the corresponding blue projection) has energy above
random.
\begin{itemize}
  \item If that energy is carried by a structured seed
        (line, function graph, vertical or horizontal packing),
        Theorem~15 of \texttt{incidences.tex} already gives
        $|I|=o(k_{\star})$.
  \item If that energy is carried by an exact cylinder or a medium
        ($|B_x|\ge q/3$) approximate cylinder, the seed piece has size
        $O(q\log q)$ and the lift is $O(\sqrt{n}\,(\log n))$, below SOTA
        only after the same polylog bookkeeping as Theorem~\ref{thm:weak}.
  \item If that energy is intermediate on thin fibres, $P$ is an instance
        of the leftover in \texttt{inverse-energy.tex}.
\end{itemize}
In every case the medium-star bucket reduces to a seed statement already
isolated. It does not close Conjecture~\ref{conj:med}.
\end{proposition}

The Sidon Cauchy--Schwarz bound
$\sum_u D(u)^{2}=O(d|P|+|P|^{2})$ is still not $o(|I|^{2})$ at
$|I|\sim\Delta=\Theta(\sqrt{n\log n})$:
the main term $|P|^{2}$ saturates the target, and the $d|P|$ term is
$|I|\cdot\Delta$ after the lift, the same first-moment count as
Lemma~\ref{lem:med}.
A genuine saving would need either a smaller energy increment on $P$
(the leftover) or a second-moment restriction on how $I$ sits over $P$.

We do \emph{not} claim the medium bucket is closed.

\section{Heavy stars}

\begin{lemma}[Few heavy stars]\label{lem:heavy-few}
$|H_I|=O(\log\log n)$.
Each heavy star has size at most $\Delta=O(\sqrt{n\log n})$.
\end{lemma}

\begin{proof}
Lemma~\ref{lem:nstars}.
\end{proof}

A set $I$ that is covered by $O(\log\log n)$ stars of size $O(\sqrt{n\log n})$ can have size $O(\sqrt{n\log n}\,\log\log n)$, a $\log\log$ above the SOTA target.
To remove the $\log\log$ one uses the same convexity as~\cite[Lemma~3.4]{HHKP2025}: the integer vector $(d_I(w))_{w\in H_I}$ is majorized by a vector with $O(1)$ entries of size $\Theta(\Delta)$ (a vertex of $G_2$ has degree $\le\Delta$, and the $\ell_{1}$-mass $\sum_{H_I}d_I\le|I|\Delta$ cannot support many huge stars).
More precisely, $\sum_{w\in H_I}d_I(w)\le|I|\Delta$ and each $d_I(w)\le\Delta$, so the number of stars that are actually of size $\Theta(\Delta)$ is $O(|I|/\Delta)=O(1)$ at the SOTA scale $|I|\sim\Delta$.
Those $O(1)$ huge stars contribute $O(\Delta)=O(\sqrt{n\log n})$ vertices, which is the target scale itself --- they do not overshoot, but they saturate it.

\begin{lemma}[Heavy contribution]\label{lem:heavy-size}
The vertices of $I$ that lie in a heavy star number at most
\[
  \bigl|I\cap\bigcup_{w\in H_I}N_{G_2}(w)\bigr|
  \;\le\;
  O(\sqrt{n\log n}).
\]
\end{lemma}

\begin{proof}
At most $O(1)$ stars of size $\Theta(\Delta)$ plus $O(\log\log n)$ stars of size $o(\Delta)$ from the tail $t_{1}<\cdot\le\Delta/2$, the latter contributing $o(\sqrt{n\log n}\,\log\log n)$ which we thin by redefining $t_{1}=\Delta/\log n$ if needed.
The $O(1)$ huge stars give $O(\Delta)$.
\end{proof}

This is a theorem at the cost of a slightly worse cutoff, and it is the reason a heavy-star covering cannot by itself produce an independent set larger than the SOTA scale.

\section{Reduction to $G_2$-independence plus an error}

\begin{lemma}[Independent kernel]\label{lem:kernel}
Let $I\subset W$ and let $I_0\subset I$ be a maximal $G_2$-independent subset.
Then every vertex of $I\setminus I_0$ has a $G_2$-neighbour in $I_0$, and
\[
  |I|
  \;\le\;
  \alpha(G_2)
  \;+\;
  \bigl|I\cap\bigcup_{w\in H_I}N_{G_2}(w)\bigr|
  \;+\;
  |I_{\mathrm{closed}}|,
\]
where $I_{\mathrm{closed}}$ is the set of vertices of $I\setminus I_0$ whose every $G_2$-edge into $I_0$ is closed by a light or medium star.
\end{lemma}

\begin{proof}
$|I_0|\le\alpha(G_2)$.
Partition $I\setminus I_0$ according to whether some $G_2$-edge into $I_0$ is open (then $I$ already spans an open edge and we are done for the original lemma) or every such edge is closed.
In the latter case each vertex lies in a closed star meeting $I_0$.
Stars in $H_I$ are charged to the heavy term; the rest to $I_{\mathrm{closed}}$.
\end{proof}

\begin{theorem}[What the buckets give]\label{thm:buckets}
Assume the seed bound $\alpha(G_R)=O(q\log q)$ (Theorem~14 of \texttt{energy-increment.tex}, conditional on the inverse-energy lemma).
Then $\alpha(G_2)=O(\sqrt{n}\,(\log n))$ and, for every $I$,
\[
  |I|
  \;\le\;
  O\bigl(\sqrt{n}\,(\log n)\bigr)
  \;+\;
  O(\sqrt{n\log n})
  \;+\;
  |I_{\mathrm{closed}}|.
\]
If Conjecture~\ref{conj:med} holds and the light-star second moment saves $n^{2\eps}$, then $|I_{\mathrm{closed}}|=o(k)$ at the SOTA scale $k=\Theta(\sqrt{n\log n})$, and every $I$ of that size either spans an open $G_2$-edge or contradicts the display.
\end{theorem}

The open-edge lemma of \texttt{family-a-independence.tex} is therefore reduced to Conjecture~\ref{conj:med} plus a light-star second moment.
Both are geometric (incidences between $I$ and short parabolic chords in the $\mathrm{GL}_2$ lift), not probabilistic.

\section{Hitting leftover open pairs}

If $I$ is not $G_2$-independent and some $G_2$-edge in $I$ is open, we are done.
If $I$ \emph{is} $G_2$-independent, the incidence theorem of \texttt{incidences.tex} already gives $|I|=o(k_{\star})$ for every structured pair of seed projections, and the energy increment of \texttt{energy-increment.tex} is the remaining seed-side input.

The mixed case --- $I$ has $G_2$-edges, all closed --- is Theorem~\ref{thm:buckets}.

\section{Unconditional corollary (weaker target)}

\begin{theorem}[Explicit bound under only the seed dichotomy]\label{thm:weak}
Assume $\alpha(G_R)=O(q\log q)$ and ignore cleanup (i.e.\ bound $\alpha(G_2)$, not $\alpha(A_q)$).
Then
\[
  \alpha(G_2)\;=\;O\bigl(\sqrt{n}\,(\log n)\bigr)
\]
and the triangle-free graph $A_q$ still satisfies only $\alpha(A_q)\ge\alpha(G_2)$; the upper bound on $\alpha(A_q)$ requires the open-edge lemma.
If in addition every closed-edge error satisfies $|I_{\mathrm{closed}}|=O\bigl(\sqrt{n}\,(\log n)^{4}\bigr)$ --- which follows from Lemma~\ref{lem:light} with $t_{3}=(\log n)^{c}$ and no medium stars of size $n^{1/4}$ --- then
\[
  \alpha(A_q)\;=\;O\bigl(\sqrt{n}\,(\log n)^{4}\bigr)
\]
and
\[
  R(3,k)\;=\;\Omega\Bigl(\frac{k^{2}}{(\log k)^{8}}\Bigr)
\]
via a polynomial-time family, improving Alon's $\Omega(k^{3/2})$.
\end{theorem}

The medium-star hypothesis ``no $M_I$ of size $n^{1/4}$'' is stronger than Conjecture~\ref{conj:med} and is not proved.
Theorem~\ref{thm:weak} is the shape of the bound we get if the geometric leftovers cooperate at polylog scale rather than at $n^{1/4}$.

\section{What is not claimed}

The open-edge lemma at the SOTA scale $C\sqrt{n\log n}$ is not a theorem.
What is a theorem is the kernel reduction (Lemma~\ref{lem:kernel}), the heavy-star bound (Lemma~\ref{lem:heavy-size}), and the observation that light/medium buckets are the only remaining geometric input.
Family~L remains the only complete SOTA object.

\begin{thebibliography}{9}
\bibitem{HHKP2025}
Z.~Hefty, P.~Horn, D.~King and F.~Pfender,
\emph{Improving $R(3,k)$ in just two bites},
arXiv:2510.19718, 2025.
\end{thebibliography}

\end{document}
